\paragraph{La densidad de las líneas retiradas.}
En la primera media de \(\mathcal M_n(b)\) se suma sobre los estados de los dos cruces y
\(\sum_s q_{o,s}=1\). La ley de motivos y el escalado en la llegada regular, desarrollados
en el apéndice~\ref{app:transfer}, dan la medida de pares
\(w_{o,b}\,dr\,dc\,dh\) para cada orientación. En efecto, las cuatro etiquetas axiales
aportan \(b^4\), cada diagonal distinta aporta una vacancia y los dos factores \(jD_z\)
convergen a \(W_b^0(z)\) en el soporte. El factor exterior \(n/j^4\) y los \(n^3\) de la
suma de Riemann dejan \(n^4 b^4/j^4\to1\). Por simetría en los extremos del par, sobre
cada línea de coordenada \(t\) su masa se integra como
\(\frac12\iint p_b(z)p_b(z')\,dr\,dr'/U_o(t)\). Al sumar las líneas resulta \(W_b\).

Si cada par contribuye una función de prueba \(a\) en cada una de sus dos filas,
la simetría de \(a(r)+a(r')\) cancela ese factor \(1/2\) al integrar un solo extremo.
En la familia \(o\), integrar el otro extremo da \(m_{o,b}(t)/U_o(t)\). Así,
\[
 \sum_o\int_\Omega w_{o,b}\{a(r_e)+a(r_f)\}\,dr\,dc\,dh
 =\iint a(r)p_b(r,c)\gamma_b(r,c)\,dr\,dc
 =\int a(r)\theta_b(r)\,dr .
\]
Por transposición se obtiene la misma densidad para las columnas. En particular,
\(\int\theta_b=2W_b\). La medida de filas perdida por los pares tiene densidad
\(\theta_b\), y la renormalización de las filas supervivientes repone la masa \(2W_b\)
con densidad constante. Ambas operaciones dan la dirección \(\eta_b=2W_b-\theta_b\).

El centrado se puede verificar exactamente en una malla determinista de \(m\) puntos por
lado. Sean \(p_m\) el núcleo de \(v_b\) escalado con la medida uniforme \(A_m\) de masa
\(1/m\) por punto y \(W_m,\theta_m\) las sumas de malla que corresponden a \(W_b,\theta_b\).
En una línea \(\ell\) de la familia \(o\) damos a cada par ordenado el peso
\(\omega=p_m(e)p_m(f)/(2m^3 U_o(\ell))\), incluidos los pares \(e=f\) en cada familia.
Sumar sobre ambos extremos da \(\sum\omega=W_m\). Fijar la fila \(i\) del primero y
multiplicar por dos da \(\theta_m(i)/m\). Con \(h_R^m,h_C^m\) las influencias centradas
en esta malla y \(H_m(\parmar)\) su suma en los cuatro extremos, se obtiene por linealidad
\[
 \sum\omega H_m(\parmar)
 =\frac1m\sum_i\theta_m(i)(h_R^m(i)+h_C^m(i))
 =-D\mathfrak F_m[\eta_m A_m,\eta_m A_m],\qquad \eta_m=2W_m-\theta_m .
\]
La igualdad vale para las primeras variaciones. Los pares \(e=f\) de esta cuadratura
tienen masa total \(O_b(m^{-1})\). Los menores físicos siempre usan dos casillas distintas.

\begin{proposicion}[Primera media a corte fijo]\label{prop:pt-corte}
Manténgase \(v_b\) fijo y evalúese \(\mathfrak F\) con las dos medidas
\((1+\varepsilon\eta_b)\,dr\), \((1+\varepsilon\eta_b)\,dc\), reescalando \(v_b\) en
cada \(\varepsilon\). Denótese por \(dF_2(b)\) la derivada en cero. Entonces
\begin{align}
 \lim_{n\to\infty}\frac n{j^4}\mathbb E_\mu
 \Big[\sum_{\parmar} w_{\parmar}\,j\Delta_{\parmar}\mathfrak F\,;F_b\Big]&=dF_2(b),
 \label{eq:pt-menores-limite}\\
 P_T(b):=\lim_{n\to\infty}-n\mathbb E_\mu
 \Big[\sum_{\parmar}(\pi_{\parmar}^{\rm ex}-\pi_{\parmar}^T)\,;F_b\Big]&=(1/6-2F(b))W_b-dF_2(b).
 \label{eq:pt-frontera-fija}
\end{align}
\end{proposicion}
\begin{proof}
Para reemplazar \(H_M(\parmar)\) por las influencias del perfil hay que transportar también la
derivada del escalado. En las series del funcional, el núcleo centrado \(\centrado=jD-1\) tiene
brecha \(\|\centrado\|_{2\to2}\le r_b<1\) y secciones de norma a lo sumo \(C_b\). Un recorrido
de longitud \(k\ge2\) entre extremos fijos queda acotado por \(C_b^2r_b^{k-2}\).
Con \(T_R=(I-\centrado\centrado^*)^{-1}-I\), en la diagonal de \(T_R-\centrado\centrado^*\) se ha
 sustraído la única
vuelta de longitud dos. En columnas se intercambian \(\centrado\) y \(\centrado^*\).
Los términos truncados del funcional son por ello sumas de grafos simples.
Derivar una medida fija un vértice y resta su media. Derivar el escalado en una fila
\(i\) añade un camino hacia esa raíz, pues la respuesta de columna en el calibre centrado es
\(-(I-\centrado^*\centrado)^{-1}\centrado_{i,\cdot}\). Las elecciones de vértice o arista derivada
 multiplican
la cola geométrica por un polinomio en la longitud, por lo que sigue siendo sumable.

A truncación fija las coincidencias accidentales de etiquetas cuestan \(O_{b,k}(1/n)\).
La ley de motivos en las líneas vivas, uniforme en la raíz, lleva cada suma al mismo grafo
con núcleo \(p_b\). La excepción a precisión fija en \(F_b\) tiene masa exponencial y los
pesos escalados están acotados. Se toma \(n\to\infty\) y después se retira la truncación.
Así se reemplazan en la media ponderada \(\mathfrak F(M),h_R^M,h_C^M\) por \(F(b),h_R,h_C\),
con error \(o_b(1)\). Aquí \(h_R,h_C\) son las influencias continuas, de media cero.
Al integrar los cuatro extremos, su contribución es
\[
 \int h_R\theta_b\,dr+\int h_C\theta_b\,dc
 =-D\mathfrak F[\eta_b\,dr,\eta_b\,dc]=-dF_2(b).
\]
Esto prueba~\eqref{eq:pt-menores-limite} por \(j\Delta_{\parmar}\mathfrak F=-H_M(\parmar)+O(j^{-1})\).
Sustituirlo en~\eqref{eq:pt-par-fisico} y usar la masa \(W_b\)
prueba~\eqref{eq:pt-frontera-fija}.
\end{proof}
La identidad~\eqref{eq:pt-frontera-fija} vale para el límite de la llegada física en
\(F_b\). Si se fijan las cuatro etiquetas axiales y se exigen las tres diagonales vacías
en la solución completa, la comparación de ambas medias cuesta \(O(b)+o_b(1)\).
La positividad uniforme de \(v_b\) y la continuidad del escalado y de su primera variación
dan \(P_T=\lim_{b\downarrow0}P_T(b)\). En este término se varían las medidas de línea.
En \(F'(0)\), que interviene en el defecto interior, se deriva el núcleo \(v_b\) a lo
largo del perfil con ambas medidas fijas.
